\begin{proposicion}[Cota por soporte y por número de ciclos]\label{prop:soportes-llegada}
Sean \(U,v,c_s,d_i\) como en los Lemas~\ref{lem:rango-llegada} y~\ref{lem:tipos-llegada}.
Sea \(\widehat e_U\) el logaritmo de la suma de los pesos de la
Proposición~\ref{prop:menores-llegada} en cada cara, menos el cambio \(\Lambda\)
del predictor al activar los pares de esa cara, con la misma normalización basal.
Supóngase que cada sitio tiene un vecino diagonal en \(U\). Para \(v\le8\log n\),
\[
 |\Delta_U\widehat e_U|\le C^v(Cv)^{c_s+1}j^{-p}\prod_{i\in U}d_i,\qquad
 p=\begin{cases}v+c_s+1&\text{bosque},\\v+c_s&\text{un ciclo},\\
 v+c_s-1&\text{dos o más ciclos independientes}.\end{cases}
\]
Para \(v\le4\) los mismos órdenes valen para el error físico \(e\).
\end{proposicion}
\begin{proof}
Introducimos una variable \(s\) de tamaño y las actividades independientes \(\zeta_i\),
tras extraer el campo afín común de las caras. En el valor físico \(s=1/j\) y
\(|\zeta_i|\le C/j\). Una fibra de tamaño \(D_f\) cumple
\(\sum_{i\in f}d_i^{-1}\le1\). Con \(\|z\|_*=\max_i d_i|z_i|\), el cambio de carga
de las fibras es \(\le4\|z\|_*\) por vecindad y las medidas del balanceo exterior
varían en \(O(v|s|)\). Su inversa centrada y los propagadores quedan acotados como en
la Proposición~\ref{prop:menores-llegada}. En \(|s|<\kappa/v,\ |\zeta_i|<\kappa'/d_i\),
la ecuación local de las odds tiene derivada en norma \(*\) a lo sumo \(C\kappa'\).
La Hessiana normalizada del prefactor es \(I+E\), \(\|E\|_\infty\le C\kappa'\);
la serie de trazas acota su logaritmo por \(v\sum_{k\ge1}\|E\|_\infty^k/k\).
En el conteo aproximado se incorpora cada sitio por la identidad vacío más ocupado:
el cociente ocupado/vacío lleva el factor \(\zeta_i\); al prohibir sus vecinos, las
variaciones suman \(O(\sum_{l\sim i}|\zeta_l|)=O(\kappa')\). En un átomo las supresiones
entran en la variable de medidas \(\theta=sT_U\), de radio fijo en variación total.
Las derivadas de orden \(h\) del campo exterior, con su polo \(1/s\), dan diferencias
mixtas \(\le h! C^h|s|^{h-1}\) al integrar en las \(h\) direcciones suprimidas.
Sumarlas sobre subconjuntos de \(h-1\) sitios adicionales cuesta una serie geométrica
con factor polinómico \(h\) si \(Cv|s|<1\). Si la razón ocupado/vacío es a lo sumo
\(K|\zeta_i|<1/2\) en volúmenes menores, las marginales complejas sucesivas tienen
módulo a lo sumo uno. Su producto controla cada indicador de sitios prescritos. El
cambio del campo exterior al ocupar \(i\) cuesta a lo sumo
\(\prod_D(1+|e^{a_D}-1|)\le e^{\sum_D|a_D|}\), con \(a_D\) las diferencias mixtas.
Al retirar los vecinos se paga \(\exp(2K\sum_{l\sim i}|\zeta_l|)\). Fijando primero
\(K\) para el campo exterior y después \(\kappa'\) pequeño se cierra la inducción.
Así los cocientes de particiones son \(1+O(|\zeta_i|)\) y el logaritmo es \(O(v)\)
por cara. Los factores de silla y cofactor coinciden en dos caras cuando \(\zeta_i=0\).
Por tanto el alternado es holomorfo, acotado por \(C^v\) y divisible por \(\prod_i\zeta_i\).

Queda contar los grados en \(s\). Al suprimir las fibras marcadas, una masa adicional
de componente entra como \(s y_f\), donde \(y_f\) suma las ocupaciones de esa fibra.
En la acción hay una potencia exterior \(s^{-1}\); en el prefactor la potencia inicial
es \(s^0\). Así una interacción ordinaria de coste \(\ell\) tiene a lo sumo
\(\ell+1\) factores de media en la acción y \(\ell\) en el prefactor. El término que
usa a la vez las dos medidas marcadas en el segundo momento del núcleo tiene amplitud
común \(s^2 z^2/2\). Para las masas \(y_i,y_h\), escribamos
\(F(y)=(1-y)^2-1,\ L(y)=-y\). Sus polinomios son \(F_iF_h\) en la acción y
\(F_iF_h-L_iF_h-F_iL_h\) en el prefactor. Su multiplicador en coste \(\ell\)
tiene grado \(\le\ell-2\). Los términos de exceso positivo ya llevan \(s\), por lo
que pueden tocar a lo sumo \(\ell-1\) grupos en ese coste.

El logaritmo factoriza las contribuciones sobre grupos separados en una suma. Para un
monomio conexo de coste total \(w\), cuya interacción \(a\) cuesta \(\ell_a\) y toca
\(h_a\) grupos, la fórmula de Euler da el número de grupos tocados,
\[
 1+w-\sum_a(\ell_a+1-h_a)-L_I ,
\]
donde \(L_I\) cuenta los ciclos de su grafo bipartito de incidencias. Cada déficit es
no negativo; al faltar un grupo la proyección se anula. El primer coste es al menos \(c_s-1\).

En un bosque las ramas de los factores de cupo uno se marginalizan exactamente. Centremos
los sitios en la media de la silla inclinada. Con \(p_0\) patas centradas, \(u\) factores
de media, \(h\) bloques de cumulantes y \(n_R,n_A\) vértices ordinarios de acción y
prefactor, hay \(L=p_0-n_R-n_A-h+1\) ciclos. Como \(p_0+u\le w+n_R\), se tocan
a lo sumo \(u+h\le w+1-n_A-L\) grupos. Los bloques centrados y los vértices de acción
tienen grado al menos dos. Con un solo vértice de prefactor no cabe un árbol con una
sola hoja posible. Si \(n_A=0,L=1\), todos los grados son dos. La serie de esos ciclos
es \(-\tfrac12\log\det(I-VR'')\), con \(V\) la covarianza por grupos; el prefactor
añade el mismo medio logaritmo con signo opuesto. Las decoraciones del campo de enlace
de coste \(t\) añaden a lo sumo \(t\) grupos, por derivar la acción de grado \(\le t+1\).
Con una interacción excepcional de déficit uno, tocar \(w\) grupos exige un árbol:
al podar una rama, las medias de sus hojas multiplican la derivada en el campo del
enlace, tanto en el conteo como en la silla; los bloques cruzados de su Hessiana tienen
traza nula. En la raíz las dos fibras están en grupos distintos. Para \(m_i=\mathbb Ey_i\),
\(v_i=m_i(1-m_i)\), la contribución estacionaria tras la traza es
\(2F_iF_h+(m_i+v_i)F_h+(m_h+v_h)F_i=0\), pues \(F(m_i)=-m_i-v_i\).
En los átomos binarios la suma de ambos polinomios vale cero. La identidad persiste al
derivar en los campos. Así el bosque toca como máximo \(w-1\) grupos.

En el caso unicíclico, al coste crítico \(c_s-1\) se podan las ramas exteriores. En el
ciclo las medias binarias \(x_i\) dan la transición de \(i\) a su sucesor \(h\),
\(\left(\begin{smallmatrix}1-t&t\\1&0\end{smallmatrix}\right)\), \(t=x_h/(1-x_i)\).
La traza del producto es \(1+\theta\), con
\(\theta=(-1)^{|\mathcal C|}\prod_{i\in\mathcal C}x_i/(1-x_i)\).
Para el cofactor sea \(E\) la matriz que copia sitios en incidencias y \(C_f\) la
covarianza categórica en la fibra dividida por su diagonal de varianzas, por filas.
Con \(C=\bigoplus_f C_f\), el lema del determinante transforma el cofactor local
\(\det C\,\det[I-E^{\mathsf T}E+E^{\mathsf T}C^{-1}E]\) en
\(\det[I-(C-I)(EE^{\mathsf T}-I)]\). Una vuelta cambia de sitio y de fibra en
cada paso, así que sólo sobreviven las dos orientaciones del ciclo, cada una de peso
\(\theta\). El determinante vale \((1-\theta)^2\). Queda en el error
\(\log(1-\theta^2)\); \(x_i\) es divisible por \(\zeta_i\), y
\(\zeta_i\partial_{\zeta_i}\) conserva esa multiplicidad doble.

Estas identidades de coeficientes persisten al truncar los excesos en \(q=4v+12\):
la recurrencia de los momentos da orden \(\lceil k/2\rceil\) para \(\mu_k\), y aun
después del prefactor \(s^{-1}\) el cambio empieza en orden \(\ge q-1>c_s+1\).
La identidad holomorfa prolonga las anulaciones a los radios indicados. Si \(R=\kappa/(2v)\)
y \(r_i=\kappa'/(2d_i)\), Cauchy en \(s\) y Schwarz en cada campo acotan la cola desde
\(w\) por \(C^v\prod_i(|\zeta_i|/r_i)(|s|/R)^w/(1-|s|/R)\).
Los costes iniciales son \(c_s+1,c_s,c_s-1\) respectivamente, salvo el coeficiente
unicíclico de coste \(c_s-1\), que gana \(\prod_{i\in\mathcal C}|\zeta_i|/r_i
\le Cv^2/j\). En el punto físico se obtienen los exponentes enunciados. Para \(v\le4\)
la comparación real por cara cuesta \(2^v\eta_v=O(j^{-p})\), también a tamaño entero.
\end{proof}
